\chapter{Wilson の定理}
\label{chap:wilson}

\section{概要}

本 blueprint では Wilson の定理の形式的証明をトップダウン方式で構築します。

\begin{theorem}[Wilson の定理]
\label{thm:wilsons_theorem}
\lean{WilsonProof.wilsons_theorem}
\uses{thm:wilsons_lemma, thm:prime_of_factorial}
素数 $p \geq 2$ に対して、次が成り立つ:
\[
  p \text{ が素数} \iff (p-1)! \equiv -1 \pmod{p}
\]

\end{theorem}

\begin{proof}
\uses{thm:wilsons_lemma, thm:prime_of_factorial}

前件と後件の双方向から示す。
\end{proof}

\section{前件: 素数ならば Wilson 合同式}

\begin{theorem}[Wilson の補題]
\label{thm:wilsons_lemma}
\lean{WilsonProof.wilsons_lemma}
\uses{lem:factorial_cast_eq_prod, lem:prod_nonzero_eq_neg_one}
素数 $p$ に対して、$(p-1)! \equiv -1 \pmod{p}$。

\end{theorem}

\begin{proof}
\uses{lem:factorial_cast_eq_prod, lem:prod_nonzero_eq_neg_one}

補題 \ref{lem:factorial_cast_eq_prod} と \ref{lem:prod_nonzero_eq_neg_one} を組み合わせると、
$(p-1)!$ を $\Zmod{p}$ に cast したものは $-1$ に等しい。
\end{proof}

\begin{lemma}
\label{lem:factorial_cast_eq_prod}
\lean{WilsonProof.factorial_cast_eq_prod}
\uses{}
素数 $p$ に対して、$(p-1)!$ の $\Zmod{p}$ への cast は
$\Zmod{p}$ の全非零元の積に等しい。
すなわち、
\[
  (p-1)! \equiv \prod_{x \in (\Z/p\Z)^*} x \pmod{p}
\]

\end{lemma}

\begin{proof}

$\{1, 2, \ldots, p-1\}$ は $(\Z/p\Z)^* = (\Z/p\Z) \setminus \{0\}$ の完全代表系をなす。
積の定義から直接従う。
\end{proof}

\begin{lemma}
\label{lem:prod_nonzero_eq_neg_one}
\lean{WilsonProof.prod_nonzero_eq_neg_one}
\uses{lem:sq_eq_one_iff_pm_one}
素数 $p$ に対して、$\Zmod{p}$ の全非零元の積は $-1$ に等しい。
\[
  \prod_{x \in (\Z/p\Z)^*} x = -1
\]

\end{lemma}

\begin{proof}
\uses{lem:sq_eq_one_iff_pm_one}

$\Zmod{p}$ の各非零元 $x$ はその逆元 $x^{-1}$ とペアを組む。
$x \neq x^{-1}$（すなわち $x^2 \neq 1$）の場合、$x$ と $x^{-1}$ の積は $1$ なので積への寄与は $1$。
$x = x^{-1}$（すなわち $x^2 = 1$）の場合、補題 \ref{lem:sq_eq_one_iff_pm_one} より
$x = 1$ または $x = -1$。
したがって全体の積は $1 \times (-1) = -1$。
\end{proof}

\begin{lemma}
\label{lem:sq_eq_one_iff_pm_one}
\lean{WilsonProof.sq_eq_one_iff_pm_one}
\uses{}
素数 $p$ と $\Zmod{p}$ の元 $x$ に対して、$x^2 = 1$ であることと
$x = 1$ または $x = -1$ であることは同値。

\end{lemma}

\begin{proof}

$p$ が素数のとき $\Zmod{p}$ は体。
$x^2 - 1 = (x-1)(x+1) = 0$ より $x = 1$ または $x = -1$（体では零因子が存在しない）。
\end{proof}

\section{後件: Wilson 合同式ならば素数}

\begin{theorem}
\label{thm:prime_of_factorial}
\lean{WilsonProof.prime_of_factorial}
\uses{}
$p \geq 2$ かつ $(p-1)! \equiv -1 \pmod{p}$ ならば $p$ は素数である。

\end{theorem}

\begin{proof}

対偶を示す。$p$ が合成数であるとする。
ある $1 < d < p$ が存在して $d \mid p$。
このとき $d \leq p - 1$ なので $d \mid (p-1)!$。
よって $(p-1)! \equiv 0 \pmod{d}$。
もし $(p-1)! \equiv -1 \pmod{p}$ ならば $-1 \equiv 0 \pmod{d}$、
すなわち $d \mid 1$ となり $d = 1$ で矛盾。
\end{proof}

%%%%%%%%%%%%%%%%%%%%%%%%%%%%%%%%%%%%%%%%%%%%%%%%%%%%%%%%%%%%%%%%%%%%%%%%
\chapter{有限生成アーベル群の構造定理}
\label{chap:abelian_group}

\section{概要}

本章では有限生成アーベル群の構造定理の形式的証明を構築します。

\begin{theorem}[有限生成アーベル群の構造定理]
\label{thm:structure_theorem_fg_abelian}
\lean{AbelianGroupStructure.structure_theorem_fg_abelian}
\uses{lem:int_is_pid, thm:structure_theorem_pid}
任意の有限生成アーベル群 $G$ は次の形に分解される:
\[
  G \cong \Z^n \oplus \bigoplus_{i} \Z / p_i^{e_i}\Z
\]
ここで $n \geq 0$ は自由部分のランク、$p_i$ は素数、$e_i$ は正整数。
\end{theorem}

\begin{proof}
\uses{lem:int_is_pid, thm:structure_theorem_pid}
$G$ は有限生成アーベル群であるから、$\Z$-加群とみなせる。
$\Z$ は PID であるから（補題 \ref{lem:int_is_pid}）、
PID 上の有限生成加群の構造定理（定理 \ref{thm:structure_theorem_pid}）を適用すると、
$G \cong \Z^n \oplus \bigoplus_i \Z / p_i^{e_i}\Z$ が得られる。
\end{proof}

\section{基礎補題}

\begin{lemma}[$\Z$ は PID]
\label{lem:int_is_pid}
\lean{AbelianGroupStructure.int_is_pid}
\uses{}
整数環 $\Z$ は単項イデアル整域（PID）である。
\end{lemma}

\begin{proof}
$\Z$ はユークリッド整域であり、ユークリッド整域は PID である。
\end{proof}

\begin{lemma}[有限生成ねじれアーベル群の有限性]
\label{lem:fg_torsion_is_finite}
\lean{AbelianGroupStructure.fg_torsion_is_finite}
\uses{}
有限生成かつねじれアーベル群は有限群である。
\end{lemma}

\begin{proof}
有限生成 $\Z$-加群がねじれ加群であるとき、自由部分のランクは $0$ である。
したがって全体が有限個の巡回群の直和に分解され、有限群となる。
\end{proof}

\section{ねじれ加群の分解}

\begin{theorem}[ねじれ加群の素冪分解]
\label{thm:torsion_module_decomposition}
\lean{AbelianGroupStructure.torsion_module_decomposition}
\uses{lem:int_is_pid}
PID $R$ 上の有限生成ねじれ加群 $M$ は、既約元の冪によるイデアル商の直和に同型である:
\[
  M \cong \bigoplus_{i} R / (p_i^{e_i})
\]
ここで各 $p_i$ は $R$ の既約元。
$R = \Z$ の場合、これは「ねじれアーベル群は素冪位数の巡回群の直和に分解される」ことを意味する。
\end{theorem}

\begin{proof}
\uses{lem:int_is_pid}
PID 上のねじれ加群に対する主イデアル分解を適用する。
まず加群の零化イデアルを素イデアル冪に因数分解し、
対応する素冪ねじれ部分加群の内部直和として表す。
各素冪ねじれ部分加群は巡回加群の直和に分解される。
\end{proof}

\section{PID 上の構造定理}

\begin{theorem}[PID 上の有限生成加群の構造定理]
\label{thm:structure_theorem_pid}
\lean{AbelianGroupStructure.structure_theorem_pid}
\uses{thm:torsion_module_decomposition}
PID $R$ 上の有限生成加群 $M$ は次の形に分解される:
\[
  M \cong R^n \oplus \bigoplus_{i} R / (p_i^{e_i})
\]
ここで $n$ は自由部分のランク、各 $p_i$ は $R$ の既約元。
\end{theorem}

\begin{proof}
\uses{thm:torsion_module_decomposition}
$M$ のねじれ部分加群 $T(M)$ を取り出すと、$M / T(M)$ はねじれ自由な有限生成 $R$-加群であり、
PID 上のねじれ自由有限生成加群は自由加群であるから $M / T(M) \cong R^n$。
射影加群の持ち上げにより $M \cong R^n \oplus T(M)$。
$T(M)$ に対して定理 \ref{thm:torsion_module_decomposition} を適用すると結論が得られる。
\end{proof}

\section{有限アーベル群の場合}

\begin{theorem}[有限アーベル群の構造定理]
\label{thm:structure_theorem_finite_abelian}
\lean{AbelianGroupStructure.structure_theorem_finite_abelian}
\uses{thm:structure_theorem_fg_abelian}
任意の有限アーベル群 $G$ は、素冪位数の巡回群の直和に同型である:
\[
  G \cong \bigoplus_{i} \Z / p_i^{e_i}\Z
\]
\end{theorem}

\begin{proof}
\uses{thm:structure_theorem_fg_abelian}
有限アーベル群は有限生成である。
定理 \ref{thm:structure_theorem_fg_abelian} を適用すると
$G \cong \Z^n \oplus \bigoplus_i \Z/p_i^{e_i}\Z$ を得る。
$G$ が有限であるから $n = 0$（$\Z$ の直和因子があれば無限群になる）。
したがって $G \cong \bigoplus_i \Z/p_i^{e_i}\Z$。
\end{proof}

\begin{theorem}[有限アーベル群の構造定理（自然数版）]
\label{thm:structure_theorem_finite_abelian_nat}
\lean{AbelianGroupStructure.structure_theorem_finite_abelian'}
\uses{thm:structure_theorem_finite_abelian}
任意の有限アーベル群 $G$ は、位数が $2$ 以上の巡回群の直和に同型である:
\[
  G \cong \bigoplus_{i} \Z / n_i\Z \quad (\forall i,\; n_i > 1)
\]
\end{theorem}

\begin{proof}
\uses{thm:structure_theorem_finite_abelian}
定理 \ref{thm:structure_theorem_finite_abelian} の各因子 $\Z/p_i^{e_i}\Z$ において、
$p_i$ が素数かつ $e_i \geq 1$ であるから $p_i^{e_i} \geq 2$。
$n_i = p_i^{e_i}$ とおけば条件を満たす。
\end{proof}
